\paragraph{Conditional Moran image.}
The two previously reported pectinate Jacobians used different finite-population
parameters. The historical exploratory panel used $M=4$ and gives
$\det Df=-6.4949035709\times10^{-4}$, whereas the common audited setting
$M=10$ gives $\det Df=-2.6196000042\times10^{-4}$. The corresponding quartet
vectors also differ. This is a parameter mismatch, not a disagreement between
the exact recursion and its derivative.

For either rooted shape, the exact finite-$M$ conditional map is smooth in
positive branch durations. The analytic matrix-exponential derivative and
quotient rule give a nonzero numerical Jacobian at the stated $M=10$ point.
A formally validated enclosure of that determinant would, by the inverse
function theorem, prove that the conditional image contains a nonempty open
subset of the quartet simplex. The present calculations establish the
derivative and its numerical stability, but do not yet provide an
interval-arithmetic enclosure. Consequently global simplex coverage and a
rigorous positive-length overlap with a hybridization line remain open.
